% ---- pg-543 （书页 531） ----

\section{分层聚类方法}
\phantomsection\label{sec:12.3}

即使使用最大最快的计算机, 我们也几乎不可能考察所有可能的分组情况. 由于这个原因, 人们提出各种各样的聚类算法, 以期能找到“合理”的聚类而又不必考察所有可能的结构.

分层聚类方法是通过一系列或者是相继的合并, 或者是相继的分割来进行的. 聚集分层法从单个对象开始. 这样在开始时有多个对象就有多少个聚类. 将那些最相似的对象首先分组, 然后将组与组根据它们之间的相似性进行合并. 最后随着相似性不断下降, 所有的组渐渐融合为一个聚类.

分割分层法的工作方式与聚集分层法正相反. 最初由所有对象组成的一个组开始, 将它分割成两个子组, 使一个子组的对象“远离”另一子组的对象. 然后将这两个组进一步分割成不相似的组. 这一过程一直到每个对象单独成为一组时为止.

聚集和分割这两种方法的结果可以用二维空间的图来表示, 这种图称为“树图”. 我们将看到, 用树图可以说明相继完成的历次合并或分割.

在本节中, 我们将集中讨论聚集分层方法, 特别是连接法. 在文献[3]和[8]中可找到关于分割分层法和其他聚集分层法的出色的初步介绍.

连接法既适用于对项目的聚类, 也适用于对变量的聚类. 不过这种情况并非对所有聚集分层法都成立. 我们将依次讨论单连接(最小距离或最近邻点)、完全连接(最大距离或最远邻点), 以及平均连接(平均距离). \hyperref[fig:12.2]{图 12.2} 是在这三种连接准则下聚类合并的示意图.

\begin{mvfigure}
\begin{tikzpicture}[every node/.style={outer sep=1pt}]
% ---------------- (a) ----------------
\begin{scope}
\draw[dashed] (0,0) ellipse (1.15 and 0.85);
\draw[dashed] (3.0,0) ellipse (1.05 and 0.95);
\fill (-0.50,0.20) circle (1.6pt); \node[above left] at (-0.50,0.20) {1};
\fill (-0.45,-0.45) circle (1.6pt); \node[below left] at (-0.45,-0.45) {2};
\fill (3.25,0.55) circle (1.6pt); \node[above left] at (3.25,0.55) {3};
\fill (2.80,0.05) circle (1.6pt); \node[above left] at (2.80,0.05) {4};
\fill (3.40,-0.50) circle (1.6pt); \node[below left] at (3.40,-0.50) {5};
\draw (-0.45,-0.45) -- (2.80,0.05);
\node at (5.60,-0.05) {$\dfrac{\text{类间距离}}{d_{24}}$};
\node at (1.45,-1.80) {(a) 单连接};
\end{scope}
% ---------------- (b) ----------------
\begin{scope}[yshift=-3.5cm]
\draw[dashed] (0,0) ellipse (1.15 and 0.85);
\draw[dashed] (3.0,0) ellipse (1.05 and 0.95);
\fill (-0.50,0.20) circle (1.6pt); \node[above left] at (-0.50,0.20) {1};
\fill (-0.45,-0.45) circle (1.6pt); \node[below left] at (-0.45,-0.45) {2};
\fill (3.25,0.55) circle (1.6pt); \node[above left] at (3.25,0.55) {3};
\fill (2.80,0.05) circle (1.6pt); \node[above left] at (2.80,0.05) {4};
\fill (3.40,-0.50) circle (1.6pt); \node[below left] at (3.40,-0.50) {5};
\draw (-0.50,0.20) -- (3.40,-0.50);
\node at (5.60,-0.05) {$d_{15}$};
\node at (1.45,-1.80) {(b) 完全连接};
\end{scope}
% ---------------- (c) ----------------
\begin{scope}[yshift=-7.0cm]
\draw[dashed] (0,0) ellipse (1.15 and 0.85);
\draw[dashed] (3.0,0) ellipse (1.05 and 0.95);
\fill (-0.50,0.20) circle (1.6pt); \node[above left] at (-0.50,0.20) {1};
\fill (-0.45,-0.45) circle (1.6pt); \node[below left] at (-0.45,-0.45) {2};
\fill (3.25,0.55) circle (1.6pt); \node[above left] at (3.25,0.55) {3};
\fill (2.80,0.05) circle (1.6pt); \node[above left] at (2.80,0.05) {4};
\fill (3.40,-0.50) circle (1.6pt); \node[below left] at (3.40,-0.50) {5};
\draw (-0.50,0.20) -- (3.25,0.55);
\draw (-0.50,0.20) -- (2.80,0.05);
\draw (-0.50,0.20) -- (3.40,-0.50);
\draw (-0.45,-0.45) -- (3.25,0.55);
\draw (-0.45,-0.45) -- (2.80,0.05);
\draw (-0.45,-0.45) -- (3.40,-0.50);
\node at (6.10,-0.05) {$\dfrac{d_{13}+d_{14}+d_{15}+d_{23}+d_{24}+d_{25}}{6}$};
\node at (1.45,-1.80) {(c) 平均连接};
\end{scope}
\end{tikzpicture}
\phantomsection\label{fig:12.2}\mvcap{图 12.2\quad 类间距离(相异性)}
\end{mvfigure}

从图中可以看到, 当两个组根据它们最邻近成员之间的距离融合时, 便产生单连接结果; 当两个组根据它们最远成员之间的距离融合时, 便产生完全连接; 当两个组根据它们相应位置上的各对成员之间的平均距离融合时, 便产生平均连接.

下面是将 $N$ 个对象(项目或变量)分组时聚集分层聚类算法的操作步骤:

1.\quad 从 $N$ 个类(每类只含一个元素)和 $N\times N$ 对称距离(或相似性)矩阵 $\mathbf{D}=\{d_{ik}\}$ 开始.

2.\quad 从距离阵中搜寻最邻近(最相似)的两个类. 记“最相似”的两个类为 $U$ 和 $V$, 它们之间的距离记为 $d_{UV}$.
